%
% Typography and Publishing - 2nd Project
% Ac.Year 2021/2022, Summer Semester
%
% Login xkvace00, Kvacek David
% 1st year BITP, full-time, FIT
% xkvace00@stud.fit.vutbr.cz
%
% Date: 2022-03-13
%

\documentclass[twocolumn, a4paper, 11pt]{article}
\usepackage[left=1.5cm, text={18cm, 25cm}, top=2.5cm]{geometry}
\usepackage[czech]{babel}
\usepackage[utf8]{inputenc}
\usepackage[IL2]{fontenc}
\usepackage{hyperref}
\usepackage{times}
\usepackage{amsmath}
\usepackage{amsthm}
\usepackage{amsfonts}
\usepackage{amssymb}

\newtheorem{definition}{Definice}
\newtheorem{sentence}{Věta}

\begin{document}

\begin{titlepage}
    \begin{center}
        \Huge\textsc{Vysoké učení technické v~Brně} \\[0.1em]
        \huge\textsc{Fakulta informačních technologií} \\
        \vspace{\stretch{0.382}}
        \LARGE{Typografie a~publikování\,--\,2. projekt} \\[0.1em]
        \LARGE{Sazba dokumentů a~matematických výrazů} \\
        \vspace{\stretch{0.618}}
    \end{center}
    \Large{2022 \hfill David Kvaček (xkvace00)}
\end{titlepage}

\section*{Úvod}
V~této úloze si vyzkoušíme sazbu titulní strany, matematic\-kých vzorců, prostředí a~dalších textových struktur obvyklých pro technicky zaměřené texty (například rovnice (\ref{eq2}) nebo Definice \ref{def2} na straně \pageref{def2}).
Pro vytvoření těchto odkazů používáme příkazy \verb|\label|, \verb|\ref| a~\verb|\pageref|.

Na titulní straně je využito sázení nadpisu podle optického středu s~využitím zlatého řezu.
Tento postup byl probírán na přednášce.
Dále je na titulní straně použito odřádkování se zadanou relativní velikostí 0,4~em a~0,3~em.

\section{Matematický text}
Nejprve se podíváme na sázení matematických symbolů a~výrazů v~plynulém textu včetně sazby definic a~vět s~využitím balíku \verb|amsthm|.
Rovněž použijeme poznámku pod čarou s~použitím příkazu \verb|\footnote|.
Někdy je vhodné použít konstrukci \verb|${}$| nebo \verb|\mbox{}|, která říká, že (matematický) text nemá být zalomen.

\begin{definition}
    \emph{Nedeterministický Turingův stroj} (NTS) je šestice tvaru $M = (Q, \Sigma, \Gamma, \delta, q_0, q_F)$, kde:
    \begin{itemize}
        \item $Q$ je konečná množina \emph{vnitřních (řídicích) stavů,}
        \item $\Sigma$ je konečná množina symbolů nazývaná \emph{vstupní abeceda, $\Delta \notin \Sigma$,}
        \item $\Gamma$ je konečná množina symbolů, $\Sigma \subset \Gamma, \Delta \in \Gamma$, nazývaná \emph{pásková abeceda,}
        \item $\delta : (Q\,\backslash\,\{q_F\}) \times \Gamma\,\rightarrow\,2^{Q \times (\Gamma \cup \{L, R\})}$, kde ${L, R \notin \Gamma}$, je parciální \emph{přechodová funkce}, a
        \item $q_0 \in Q$ je \emph{počáteční stav} a~${q_F \in Q}$ je \emph{koncový stav.}
    \end{itemize}
\end{definition}

Symbol $\Delta$ značí tzv. \emph{blank} (prázdný symbol), který se vyskytuje na místech pásky, která nebyla ještě použita.

\emph{Konfigurace pásky} se skládá z~nekonečného řetězce, který reprezentuje obsah pásky, a~pozice hlavy na tomto řetězci.
Jedná se o~prvek množiny ${\{\gamma\Delta^{\omega}~|~\gamma \in \Gamma^*\}\times\mathbb{N}}$\footnote{Pro libovolnou abecedu $\Sigma$ je $\Sigma^{\omega}$ množina všech \emph{nekonečných} řetězců nad $\Sigma$, tj. nekonečných posloupností symbolů ze $\Sigma$.}. \\
\emph{Konfiguraci pásky} obvykle zapisujeme jako $\Delta xyz\underline{z}x\Delta$\dots \\ (podtržení značí pozici hlavy).
\emph{Konfigurace stroje} je pak dána stavem řízení a~konfigurací pásky.
Formálně se jedná o~prvek množiny ${Q\times\{\gamma\Delta^{\omega}~|~\gamma \in \Gamma^*\}\times\mathbb{N}}$.

\subsection{Podsekce obsahující definici a~větu}
\begin{definition}
    \label{def2}
    \emph{Řetězec} $w$ \emph{nad abecedou} $\Sigma$ \emph{je přijat NTS} $M$, jestliže $M$ při aktivaci z~počáteční konfigurace pásky $\underline{\Delta}w\Delta$\,\dots\,a~počátečního stavu $q_0$ může zastavit přechodem do koncového stavu $q_F$, tj. $(q_0, \Delta w\Delta^w, 0)\,\underset{M}{\overset{*}{\vdash}}\,(q_F, \gamma, n)$ pro nějaké $\gamma \in \Gamma^*$ a~$n \in \mathbb{N}$.
    
    Množinu $L(M)\,=\,\{w\,|\,w$ je přijat NTS $M\,\}\,\subseteq\,\Sigma^*$ nazýváme \emph{jazyk přijímaný NTS} $M$.
\end{definition}

Nyní si vyzkoušíme sazbu vět a~důkazů opět s~použitím balíku \verb|amsthm|.

\begin{sentence}
    Třída jazyků, které jsou přijímány NTS, odpovídá \emph{rekurzivně vyčíslitelným jazykům}.
\end{sentence}

\section{Rovnice}
Složitější matematické formulace sázíme mimo plynulý text. Lze umístit několik výrazů na jeden řádek, ale pak je třeba tyto vhodně oddělit, například příkazem \verb|\quad|.
$$x^2 - \sqrt[4]{y_1 * y_2^3}\quad x > y_1 \geq y_2\quad z_{z_z} \neq \alpha_1^{\alpha_2^{\alpha_3}}$$

V~rovnici (\ref{eq1}) jsou využity tři typy závorek s~různou explicitně definovanou velikostí.

\begin{eqnarray}
    x & = & \bigg\{a\,\oplus\,\Big[b\,\cdot\,\big(c\,\ominus\,d\big)\Big]\bigg\}^{4/2} \label{eq1} \\
    y & = & \lim_{\beta \to \infty} \frac{\tan^2 \beta - \sin^3 \beta}{\frac{1}{\frac{1}{\log_{42} x} + \frac{1}{2}}} \label{eq2}
\end{eqnarray}

V~této větě vidíme, jak vypadá implicitní vysázení limity $\lim_{n \to \infty} f(n)$ v~normálním odstavci textu.
Podobně je to i~s~dalšími symboly jako $\bigcup_{N \in \mathcal{M}} N$ či $\sum_{j = 0}^n x_j^2$.
S~vynucením méně úsporné sazby příkazem \verb|\limits| budou vzorce vysázeny v~podobě $\lim\limits_{n \to \infty} f(n)$ a~$\sum\limits_{j = 0}^n x_j^2$.

\section{Matice}
Pro sázení matic se velmi často používá prostředí \verb|array| a~závorky (\verb|\left|, \verb|\right|).

$$\textbf{A} = \begin{vmatrix} a_{11} & a_{12} & \dots & a_{1n} \\
a_{21} & a_{22} & \dots & a_{2n} \\
\vdots & \vdots & \ddots & \vdots \\
a_{m1} & a_{m2} & \dots & a_{mn} \end{vmatrix} = \begin{vmatrix}
t & u~\\ v~& w \end{vmatrix} = tw - uv$$

Prostředí \verb|array| lze úspěšně využít i~jinde.

$$\binom{n}{k} = \left\{ \begin{array}{ll}
\frac{n!}{k!(n - k)!} & \mbox{pro } 0 \leq k~\leq n \\
0 & \mbox{pro } k~> n \mbox{ nebo } k~< 0
\end{array} \right.$$

\end{document}

